% Intended LaTeX compiler: pdflatex
\documentclass[a4paper,12pt]{article}
\usepackage[utf8]{inputenc}
\usepackage[T1]{fontenc}
\usepackage{graphicx}
\usepackage{longtable}
\usepackage{wrapfig}
\usepackage{rotating}
\usepackage[normalem]{ulem}
\usepackage{amsmath}
\usepackage{amssymb}
\usepackage{capt-of}
\usepackage{hyperref}
\usepackage{\string~/.emacs.d/common}
\author{mahmood}
\date{\today}
\title{common lisp graphics}
\hypersetup{
 pdfauthor={mahmood},
 pdftitle={common lisp graphics},
 pdfkeywords={},
 pdfsubject={},
 pdfcreator={Emacs 28.2 (Org mode 9.5.5)}, 
 pdflang={English}}
\begin{document}

\maketitle
\begin{note}
this document and others can be found on my website\\
\url{https://mahmoodsheikh36.github.io/}
\end{note}
we need the libraries \texttt{sdl2}, \texttt{livesupport} (to be able to continue if an error occurs in the main loop) and \texttt{opengl} for rendering\\
using \href{20230224163920-common_lisp.org}{quicklisp}\\
\lstset{language=Lisp,label= ,caption= ,captionpos=b,numbers=none}
\begin{lstlisting}
(ql:quickload :sdl2)
(ql:quickload :livesupport)
(ql:quickload :cl-opengl)
\end{lstlisting}
here im implementing an abstraction layer over opengl/sdl2 to be able to render things like math functions easily\\
we need a generic way to represent objects that are renderable onto a window so we use the generic function \texttt{render} and implement it in the various classes\\
\lstset{language=Lisp,label= ,caption= ,captionpos=b,numbers=none}
\begin{lstlisting}
(defgeneric render (obj)
  (:documentation "a function to render the object in an opengl context"))

(defstruct rgb
  "rgb color"
  (r 0) (g 0) (b 0))

(defstruct point-2d ()
  x y (color (make-rgb :r 0 :b 0 :g 0)))

(defmethod render ((p point-2d))
  (gl:with-primitives :points
    (gl:vertex (point-2d-x p) (point-2d-y p))))

(defstruct line-2d
  point1 point2 (color (make-rgb :r 0 :b 0 :g 0)) (width 1))

(defmethod render ((line line-2d))
  (let ((point1 (line-2d-point1 line))
        (point2 (line-2d-point2 line))
        (width (line-2d-width line)))
    (gl:line-width width)
    (gl:with-primitives
        :lines
      (gl:vertex (point-2d-x point1) (point-2d-y point1))
      (gl:vertex (point-2d-x point2) (point-2d-y point2)))))

(defclass window ()
  ((draw-axis :initarg :draw-axis :accessor window-draw-axis)
   (renderables :initform '() :accessor window-renderables)
   (width :initarg :width :initform 750 :accessor window-width)
   (height :initarg :height :initform 750 :accessor window-height)))

(defmethod window-run ((my-window window))
  (sdl2:with-init (:everything)
    (format t "Using SDL Library Version: ~D.~D.~D~%"
            sdl2-ffi:+sdl-major-version+
            sdl2-ffi:+sdl-minor-version+
            sdl2-ffi:+sdl-patchlevel+)
    (sdl2:with-window
        (sdl-window
         :title "basic graphics"
         :flags '(:shown :opengl)
         :w (window-width my-window)
         :h (window-height my-window))
      (sdl2:with-gl-context (gl-context sdl-window)
        ;; init opengl
        (sdl2:gl-make-current sdl-window gl-context)

        ;; enter main loop
        (sdl2:with-event-loop (:method :poll)
          (:keyup (:keysym keysym)
                  (when (sdl2:scancode= (sdl2:scancode-value keysym) :scancode-q)
                    (sdl2:push-event :quit)))
          (:idle
           ()
           (livesupport:update-repl-link) ;; accept requests from slime/sly, for live programming
           (livesupport:continuable (window-render my-window))
           (sdl2:gl-swap-window sdl-window))
          (:quit () t))))))

(defmethod window-render ((win window))
  "render the renderables, using opengl"
  (gl:clear-color 255 255 255 1.0)
  (gl:clear :color-buffer)
  (loop for renderable in (window-renderables win)
        do (gl:matrix-mode :projection)
           (gl:load-identity) ;; reset matrix to default state
           ;; (gl:ortho -1 1 -1 1 -1 1) ;; normalize window space, not really needed, done by default
           (gl:viewport 0 0 (window-width win) (window-height win))
           (gl:matrix-mode :modelview)
           (with-slots (color) renderable
             (gl:color (/ (rgb-r color) 255) (/ (rgb-b color) 255) (/ (rgb-g color) 255)))
           (render renderable))
  (sdl2:delay 20))

(defmethod window-add-renderable ((win window) renderable)
  "add a renderable to renderables, a renderable is an object that has the generic function render"
  (push renderable (window-renderables win)))

(defstruct axis
  "a renderable, can hold other renderables which would be drawn locally (transformed), pos is the (point-2d) position in the window"
  (renderables '()) pos width height (color (make-rgb :r 0 :b 0 :g 0))
  ;; bounds of the axes
  (min-y -1) (max-y 1) (min-x -1) (max-x 1))

(defmethod axis-add-renderable ((a axis) renderable)
  "add a renderable to renderables, a renderable is an object that has the generic function render"
  (push renderable (axis-renderables a)))

(defmethod render ((a axis))
  (with-slots (min-y min-x max-y max-x width height) a
    ;; axes view matrix setup
    (gl:matrix-mode :projection)
    (gl:viewport (point-2d-x (axis-pos a))
                 (point-2d-y (axis-pos a))
                 width
                 height)
    ;; apply matirx to make the screenspace correct for original width/height
    (gl:ortho 0 width 0 height -1 1)
    ;; render ticks of both axes
    (let ((origin (make-point-2d
                   :x (map-num 0 0 width (map-num 0 min-x max-x 0 width) width)
                   :y (map-num 0 0 height (map-num 0 min-y max-y 0 height) height))))
      (loop for x from (point-2d-x origin) below width by 75
            do (render
                (make-line-2d
                 :point1 (make-point-2d
                          :x x :y (- (point-2d-y origin) 10))
                 :point2 (make-point-2d
                          :x x :y (- (point-2d-y origin) -10))
                 :width 3)))
      (loop for x from (point-2d-x origin) above 0 by 75
            do (render
                (make-line-2d
                 :point1 (make-point-2d
                          :x x :y (- (point-2d-y origin) 10))
                 :point2 (make-point-2d
                          :x x :y (- (point-2d-y origin) -10))
                 :width 3)))
      (loop for y from (point-2d-y origin) above 0 by 75
              do (render
                  (make-line-2d
                   :point1 (make-point-2d
                            :x (- (point-2d-x origin) 10) :y y)
                   :point2 (make-point-2d
                            :x (- (point-2d-x origin) -10) :y y)
                   :width 3)))
      (loop for y from (point-2d-y origin) below height by 75
              do (render
                  (make-line-2d
                   :point1 (make-point-2d
                            :x (- (point-2d-x origin) 10) :y y)
                   :point2 (make-point-2d
                            :x (- (point-2d-x origin) -10) :y y)
                   :width 3))))
    ;; reset matrix then apply another matrix so that the screenspace is correct for min/max-x/y
    (gl:load-identity)
    (gl:ortho min-x max-x min-y max-y -1 1)
    (gl:matrix-mode :modelview)
    ;; render axes
    (render
     (make-line-2d
      :point1 (make-point-2d
               :x min-x :y 0)
      :point2 (make-point-2d
               :x max-x :y 0)
      :width 3))
    (render
     (make-line-2d
      :point1 (make-point-2d
               :x 0 :y max-y)
      :point2 (make-point-2d
               :x 0 :y min-y)
      :width 3))
    ;; render children (renderables)
    (loop for renderable in (axis-renderables a)
          do (render renderable))))

(defstruct plot
  "lam is the lambda function to call with x (should return y)"
  lam
  ;; bounds of the plot (drawing bound of the plots)
  min-y max-y min-x max-x
  (color (make-rgb :r 0 :b 0 :g 0)))

(defun map-num (num src-min src-max dest-min dest-max)
  "(map-num 0.5 -1 1 -50 50) => 25.0"
  (/ (- (+ (* (- num src-min) (- dest-max dest-min)) (* dest-min src-max)) (* dest-min src-min)) (- src-max src-min)))

(defmethod render ((p plot))
  (let ((prev-point nil)
        (lam (plot-lam p)))
    (loop for x from -100 below 100
          do (let* ((x (map-num x -100 100 (plot-min-x p) (plot-max-x p)))
                    (new-y (funcall lam x))
                    (new-point (make-point-2d :x x :y new-y)))
               (if prev-point
                   (render (make-line-2d
                            :point1 prev-point
                            :point2 new-point))
                   (render new-point))
               (setf prev-point new-point)))))

(defstruct discrete-plot
  "a discrete plot, one that connects a finite amount of points"
  points (color (make-rgb :r 0 :b 0 :g 0)))

(defmethod render ((p discrete-plot))
  (let* ((points (discrete-plot-points p))
         (prev-point (elt points 0)))
    (loop for i from 1 below (length points)
          do (let ((new-point (elt points i)))
               (render (make-line-2d
                        :point1 prev-point
                        :point2 new-point))
               (setf prev-point new-point)))))
\end{lstlisting}
example usage:\\
\lstset{language=Lisp,label= ,caption= ,captionpos=b,numbers=none}
\begin{lstlisting}
(defun sdl2-test ()
  (defparameter *win* (make-instance 'window :draw-axis t :width 750 :height 750))
  (defparameter *axis* (make-axis
                        :min-x -1
                        :max-x 3
                        :max-y 3
                        :min-y -1
                        :pos (make-point-2d :x 100 :y 100)
                        :width 750
                        :height 750))
  (window-add-renderable *win* *axis*)
  (axis-add-renderable *axis* (make-plot :lam (lambda (x) (* x x)) :min-x -2 :max-x 1))
  ;; (axis-add-renderable *axis* (make-plot :lam (lambda (x) (sin x))))
  ;; (axis-add-renderable *axis* (make-plot :lam (lambda (x) (tan x))))
  ;; (axis-add-renderable *axis* (make-plot :lam (lambda (x) (cos x))))
  (axis-add-renderable
   *axis*
   (make-discrete-plot :points
                       (list (make-point-2d :x -1 :y 1)
                             (make-point-2d :x 0 :y 0)
                             (make-point-2d :x 2 :y 1))))
  (sb-thread:make-thread (lambda () (window-run *win*))))
\end{lstlisting}
we can create objects and add them dynamically to the window:\\
\lstset{language=Lisp,label= ,caption= ,captionpos=b,numbers=none}
\begin{lstlisting}
(window-add-renderable
 *win*
 (make-discrete-plot :points
                     (list (make-point-2d :x (/ (random 10) 10) :y (/ (random 10) 10))
                           (make-point-2d :x (/ (random 10) 10) :y (/ (random 10) 10)))))
"done"
\end{lstlisting}
\end{document}